% =====================================================================
% APÊNDICE A -- só as tabelas de parâmetros da simulação.
% Movidas de RESULTADOS/setup_simulacao.tex em 2026-11-07, a pedido da autora.
% O texto e as seções continuam nos Resultados, citando estas tabelas.
% =====================================================================

\begin{table}[!ht]
\caption{Channel parameters set in the MATLAB script that runs QuaDRiGa.}
\label{tab:param_canal}
\vspace{4pt} % a classe zera o espaço entre a caption e a tabela
\centering
\footnotesize
\renewcommand{\arraystretch}{1.1}
\setlength{\aboverulesep}{0pt}
\setlength{\belowrulesep}{0pt}
\begin{tabular}{|p{2.7cm}|p{3.2cm}|p{4.2cm}|p{3.1cm}|}
\toprule
\centering\textbf{Source} & \centering\textbf{Parameter} & \centering\textbf{Name in the MATLAB script} & \centering\textbf{Value} \tabularnewline
\midrule
\multirow{5}{=}{\raggedright \citesp{zhao2025jsac}} & \raggedright Number of cells and BSs & \raggedright \texttt{tx\_\allowbreak{}position} (one BS) & \raggedright 1 \tabularnewline
\cline{2-4}
 & \raggedright Sector where the UEs are & \raggedright \texttt{sector\_\allowbreak{}deg} & \raggedright $120^\circ$ \tabularnewline
\cline{2-4}
 & \raggedright Position of the interferers & \raggedright \texttt{n\_\allowbreak{}ue = K + r} & \raggedright Same rule as the served UEs \tabularnewline
\cline{2-4}
 & \raggedright Channel generator & \raggedright -- & \raggedright QuaDRiGa \tabularnewline
\cline{2-4}
 & \raggedright Channel constant within a PRB & \raggedright -- & \raggedright Flat channel \tabularnewline
\midrule
\multirow{5}{=}{\raggedright QuaDRiGa default \citesp{quadriga-manual}} & \raggedright Uses only the minimal version of the 3GPP model & \raggedright \texttt{use\_\allowbreak{}3GPP\_\allowbreak{}baseline} & \raggedright \texttt{false} \tabularnewline
\cline{2-4}
 & \raggedright Shape of the spatial correlation of the large-scale parameters & \raggedright \texttt{autocorrelation\_\allowbreak{}function} & \raggedright \texttt{Comb300} \tabularnewline
\cline{2-4}
 & \raggedright Adds the direct-path delay to the delays & \raggedright \texttt{use\_\allowbreak{}absolute\_\allowbreak{}delays} & \raggedright \texttt{false} \tabularnewline
\cline{2-4}
 & \raggedright Each path starts with a random phase & \raggedright \texttt{use\_\allowbreak{}random\_\allowbreak{}initial\_\allowbreak{}phase} & \raggedright \texttt{true} \tabularnewline
\cline{2-4}
 & \raggedright Samples per half wavelength along the track & \raggedright \texttt{sample\_\allowbreak{}density} & \raggedright 2.5 \tabularnewline
\midrule
\multirow{9}{=}{\raggedright Defined in this work} & \raggedright QuaDRiGa version & \raggedright \texttt{version} & \raggedright 2.8.1-0 \tabularnewline
\cline{2-4}
 & \raggedright Propagation scenario & \raggedright \texttt{scenario} & \raggedright \texttt{3GPP\_\allowbreak{}38.901\_\allowbreak{}UMa\_\allowbreak{}NLOS} (Table~\ref{tab:cenario_uma}) \tabularnewline
\cline{2-4}
 & \raggedright Bandwidth & \raggedright \texttt{bandwidth} & \raggedright 20 MHz \tabularnewline
\cline{2-4}
 & \raggedright Flat-channel extraction & \raggedright \texttt{fr(bandwidth, 1)} & \raggedright Response at one subcarrier \tabularnewline
\cline{2-4}
 & \raggedright BS array & \raggedright \texttt{tx\_\allowbreak{}array} & \raggedright ULA with $\lambda/2$ spacing \tabularnewline
\cline{2-4}
 & \raggedright Antenna elements (BS and UE) & \raggedright \texttt{qd\_\allowbreak{}arrayant('omni')} & \raggedright Omnidirectional \tabularnewline
\cline{2-4}
 & \raggedright Sector orientation & \raggedright \texttt{tx\_\allowbreak{}array}, \texttt{ang} & \raggedright Centered on the direction perpendicular to the array \tabularnewline
\cline{2-4}
 & \raggedright UE placement & \raggedright \texttt{rx\_\allowbreak{}track} & \raggedright Uniform over the sector area \tabularnewline
\cline{2-4}
 & \raggedright Realizations & \raggedright \texttt{rng(seed + rr)} & \raggedright New positions in each realization \tabularnewline
\bottomrule
\end{tabular}
\vspace{4pt}
\legend{Source: \citesp{zhao2025jsac}; \citesp{quadriga-manual} and the author.}
\end{table}

\begin{table}[!ht]
\caption{Parameters of the propagation scenario \texttt{3GPP\_38.901\_UMa\_NLOS}: valid range and values used in the simulation.}
\label{tab:cenario_uma}
\vspace{4pt} % a classe zera o espaço entre a caption e a tabela
\centering
\footnotesize
\renewcommand{\arraystretch}{1.2}
\setlength{\aboverulesep}{0pt}
\setlength{\belowrulesep}{0pt}
\begin{tabular}{|p{1.9cm}|p{2.6cm}|p{3.0cm}|p{2.6cm}|p{2.6cm}|}
\toprule
\centering\textbf{Set in} & \centering\textbf{Parameter} & \centering\textbf{Name} & \centering\textbf{Valid range of the propagation scenario} & \centering\textbf{Value used} \tabularnewline
\midrule
\multirow{4}{=}{\raggedright MATLAB script} & \raggedright Carrier frequency & \raggedright \texttt{center\_\allowbreak{}freq} & \raggedright 0.5 to 100 GHz & \raggedright 3.5 GHz (FR1) and 28 GHz (FR2) \tabularnewline
\cline{2-5}
 & \raggedright BS--UE distance & \raggedright \texttt{r\_\allowbreak{}min\_\allowbreak{}m}, \texttt{r\_\allowbreak{}max\_\allowbreak{}m} & \raggedright 10 to 5000 m & \raggedright 50 to 100 m \tabularnewline
\cline{2-5}
 & \raggedright BS height & \raggedright \texttt{bs\_\allowbreak{}height\_\allowbreak{}m} & \raggedright 25 m & \raggedright 25 m \tabularnewline
\cline{2-5}
 & \raggedright UE height & \raggedright \texttt{ue\_\allowbreak{}height\_\allowbreak{}m} & \raggedright 1.5 to 22.5 m & \raggedright 1.5 m \tabularnewline
\midrule
\multirow{3}{=}{\raggedright Scenario file (\texttt{.conf})} & \raggedright Number of clusters & \raggedright \texttt{NumClusters} & \multicolumn{2}{c|}{21} \tabularnewline
\cline{2-5}
 & \raggedright Paths per cluster & \raggedright \texttt{NumSubPaths} & \multicolumn{2}{c|}{20} \tabularnewline
\cline{2-5}
 & \raggedright Shadowing standard deviation & \raggedright \texttt{SF\_\allowbreak{}sigma} & \multicolumn{2}{c|}{6 dB} \tabularnewline
\bottomrule
\end{tabular}
\vspace{4pt}
\legend{Source: \citesp{quadriga-manual}; \citesp{zhao2025jsac} and the author.}
\end{table}

\begin{table}[!ht]
\caption{Signal and noise parameters.}
\label{tab:param_ruido}
\vspace{4pt} % a classe zera o espaço entre a caption e a tabela
\centering
\footnotesize
\renewcommand{\arraystretch}{1.2}
\setlength{\aboverulesep}{0pt}
\setlength{\belowrulesep}{0pt}
\begin{tabular}{|p{2.3cm}|p{3.4cm}|p{4.3cm}|p{3.2cm}|}
\toprule
\centering\textbf{Source} & \centering\textbf{Parameter} & \centering\textbf{Name in the Python simulator} & \centering\textbf{Value} \tabularnewline
\midrule
\multirow{4}{=}{\raggedright \citesp{zhao2025jsac}} & \raggedright Noise model & \raggedright \texttt{M5\_\allowbreak{}Noise.\allowbreak noise\_\allowbreak{}model} & \raggedright $\mathbf{n} \sim \mathcal{CN}(\mathbf{0}, \beta\overline{\mathbf{H}}\,\overline{\mathbf{H}}^H + \sigma^2\mathbf{I})$ \tabularnewline
\cline{2-4}
 & \raggedright SNR definition & \raggedright \texttt{M5\_\allowbreak{}Noise.\allowbreak snr\_\allowbreak{}definition} & \raggedright $10\log_{10}(E_s/\sigma^2)$ \tabularnewline
\cline{2-4}
 & \raggedright Mean symbol energy & \raggedright \texttt{M4\_\allowbreak{}Source.\allowbreak Es} & \raggedright $E_s = 1$ \tabularnewline
\cline{2-4}
 & \raggedright Thermal noise power & \raggedright \texttt{M5\_\allowbreak{}Noise.\allowbreak snr\_\allowbreak{}definition} & \raggedright $\sigma^2 = 10^{-\mathrm{SNR}/10}$ \tabularnewline
\midrule
\multirow{5}{=}{\raggedright Defined in this work} & \raggedright Channel input & \raggedright \texttt{M2\_\allowbreak{}Channel.\allowbreak backend} & \raggedright $\mathbf{H}$ and $\overline{\mathbf{H}}$ read from the \texttt{.npz} file of QuaDRiGa \tabularnewline
\cline{2-4}
 & \raggedright Channel normalization & \raggedright \texttt{M2\_\allowbreak{}Channel.\allowbreak normalization} & \raggedright Same mean power for every UE \tabularnewline
\cline{2-4}
 & \raggedright Constellation energy & \raggedright \texttt{M4\_\allowbreak{}Source.\allowbreak unit\_\allowbreak{}energy\_\allowbreak{}constellation} & \raggedright Unit mean \tabularnewline
\cline{2-4}
 & \raggedright Interferer power & \raggedright \texttt{M5\_\allowbreak{}Noise.\allowbreak beta\_\allowbreak{}convention} & \raggedright $\beta = \sigma^2(10^{\mathrm{IoTh}/10} - 1)/r$ \tabularnewline
\cline{2-4}
 & \raggedright Noise of the data symbols & \raggedright \texttt{M5\_\allowbreak{}Noise.\allowbreak data\_\allowbreak{}noise\_\allowbreak{}independent} & \raggedright Independent of the $N$ samples \tabularnewline
\bottomrule
\end{tabular}
\vspace{4pt}
\legend{Source: \citesp{zhao2025jsac} and the author. IoTh and modulation are in Table~\ref{tab:paineis}.}
\end{table}

\begin{table}[!ht]
\caption{Receiver parameters: partitioning, covariance estimation and equalization.}
\label{tab:param_receptor}
\vspace{4pt} % a classe zera o espaço entre a caption e a tabela
\centering
\footnotesize
\renewcommand{\arraystretch}{1.2}
\setlength{\aboverulesep}{0pt}
\setlength{\belowrulesep}{0pt}
\begin{tabular}{|p{2.3cm}|p{3.4cm}|p{4.3cm}|p{3.2cm}|}
\toprule
\centering\textbf{Source} & \centering\textbf{Parameter} & \centering\textbf{Name in the Python simulator} & \centering\textbf{Value} \tabularnewline
\midrule
\multirow{8}{=}{\raggedright \citesp{zhao2025jsac}} & \raggedright Antennas per cluster & \raggedright \texttt{M3\_\allowbreak{}Clustering.\allowbreak equal\_\allowbreak{}cluster\_\allowbreak{}size} & \raggedright $M_c = M/C$ \tabularnewline
\cline{2-4}
 & \raggedright Noise samples to estimate $\hat{\mathbf{R}}$ & \raggedright \texttt{M5\_\allowbreak{}Noise.\allowbreak n\_\allowbreak{}noise\_\allowbreak{}samples} & \raggedright $N = 192$ \tabularnewline
\cline{2-4}
 & \raggedright Covariance estimate & \raggedright \texttt{M6\_\allowbreak{}Covariance.\allowbreak estimator} & \raggedright Sample mean of the $N$ samples \tabularnewline
\cline{2-4}
 & \raggedright Channel knowledge at the BS & \raggedright \texttt{M7\_\allowbreak{}Equalizers.\allowbreak perfect\_\allowbreak{}csi} & \raggedright Perfect $\mathbf{H}$ \tabularnewline
\cline{2-4}
 & \raggedright Compared equalizers & \raggedright \texttt{M7\_\allowbreak{}Equalizers.\allowbreak curves} & \raggedright LMMSE, ZF and BCD-MMSE \tabularnewline
\cline{2-4}
 & \raggedright BCD-MMSE iterations & \raggedright \texttt{M7\_\allowbreak{}Equalizers.\allowbreak bcd\_\allowbreak{}iterations} & \raggedright 1 \tabularnewline
\cline{2-4}
 & \raggedright BCD update order & \raggedright \texttt{M7\_\allowbreak{}Equalizers.\allowbreak bcd\_\allowbreak{}update\_\allowbreak{}rule} & \raggedright Sequential, $c = 1, \dots, C$ \tabularnewline
\cline{2-4}
 & \raggedright Computation of the equalizer $\mathbf{W}$ & \raggedright \texttt{M1\_\allowbreak{}Driver.\allowbreak recompute\_\allowbreak{}W\_\allowbreak{}per\_\allowbreak{}realization} & \raggedright Once per realization and SNR point \tabularnewline
\midrule
\multirow{5}{=}{\raggedright Defined in this work} & \raggedright Antenna partitioning into clusters & \raggedright \texttt{M3\_\allowbreak{}Clustering.\allowbreak mapping} & \raggedright Contiguous \tabularnewline
\cline{2-4}
 & \raggedright $N$ in each test condition & \raggedright \texttt{M1\_\allowbreak{}Driver.\allowbreak test\_\allowbreak{}conditions} & \raggedright 192, 168, 132 or 24 (Table~\ref{tab:condicoes_teste}) \tabularnewline
\cline{2-4}
 & \raggedright Noise sample draws & \raggedright \texttt{M5\_\allowbreak{}Noise.\allowbreak resample\_\allowbreak{}noise\_\allowbreak{}per\_\allowbreak{}realization} & \raggedright New samples at each realization and SNR point \tabularnewline
\cline{2-4}
 & \raggedright Conditioning of $\hat{\mathbf{R}}$ & \raggedright \texttt{M6\_\allowbreak{}Covariance.\allowbreak shrinkage} & \raggedright Ledoit--Wolf shrinkage \tabularnewline
\cline{2-4}
 & \raggedright ZF form & \raggedright \texttt{M7\_\allowbreak{}Equalizers.\allowbreak zf\_\allowbreak{}form} & \raggedright $(\mathbf{H}^H\mathbf{H})^{-1}\mathbf{H}^H$ \tabularnewline
\bottomrule
\end{tabular}
\vspace{4pt}
\legend{Source: \citesp{zhao2025jsac} and the author. $C$ and $K$ are in Table~\ref{tab:paineis}. Names in the code: \texttt{dbp\_sim/params.py}.} % TODO citar Ledoit e Wolf (2004)
\end{table}

\begin{table}[!ht]
\caption{Evaluation parameters.}
\label{tab:param_avaliacao}
\vspace{4pt} % a classe zera o espaço entre a caption e a tabela
\centering
\footnotesize
\renewcommand{\arraystretch}{1.2}
\setlength{\aboverulesep}{0pt}
\setlength{\belowrulesep}{0pt}
\begin{tabular}{|p{2.3cm}|p{3.4cm}|p{4.3cm}|p{3.2cm}|}
\toprule
\centering\textbf{Source} & \centering\textbf{Parameter} & \centering\textbf{Name in the Python simulator} & \centering\textbf{Value} \tabularnewline
\midrule
\multirow{5}{=}{\raggedright \citesp{zhao2025jsac}} & \raggedright Channel realizations & \raggedright \texttt{M1\_\allowbreak{}Driver.\allowbreak n\_\allowbreak{}realizations} & \raggedright 100 \tabularnewline
\cline{2-4}
 & \raggedright Symbols per realization & \raggedright \texttt{M1\_\allowbreak{}Driver.\allowbreak n\_\allowbreak{}symbols} & \raggedright 480 \tabularnewline
\cline{2-4}
 & \raggedright Detection & \raggedright \texttt{M8\_\allowbreak{}Metric.\allowbreak detection} & \raggedright Nearest constellation point \tabularnewline
\cline{2-4}
 & \raggedright Metric & \raggedright \texttt{M8\_\allowbreak{}Metric.\allowbreak metric} & \raggedright Symbol error rate (SER) \tabularnewline
\cline{2-4}
 & \raggedright SNR range & \raggedright \texttt{M1\_\allowbreak{}Driver.\allowbreak snr\_\allowbreak{}db\_\allowbreak{}grid} & \raggedright $-10$ to $10$ dB, 2 dB step \tabularnewline
\midrule
\multirow{2}{=}{\raggedright Defined in this work} & \raggedright Bias removal before detection & \raggedright \texttt{M8\_\allowbreak{}Metric.\allowbreak bias\_\allowbreak{}removal} & \raggedright No \tabularnewline
\cline{2-4}
 & \raggedright Random seed & \raggedright \texttt{M1\_\allowbreak{}Driver.\allowbreak rng\_\allowbreak{}seed} & \raggedright 0 \tabularnewline
\bottomrule
\end{tabular}
\vspace{4pt}
\legend{Source: \citesp{zhao2025jsac} and the author. Names in the code: \texttt{dbp\_sim/params.py}.}
\end{table}
